% Elsevier Engineering Journal Template
% Suitable for: Engineering Structures, Applied Energy, Energy and Buildings,
%               Building and Environment, Construction and Building Materials,
%               Automation in Construction, Renewable Energy, Journal of Cleaner Production,
%               Computer Methods in Applied Mechanics and Engineering,
%               Mechanical Systems and Signal Processing, Structural Safety, etc.
% Last updated: 2024
%
% Uses Elsevier's official `elsarticle` document class (available on CTAN).
% Citation style: numbered [1][2] using \cite{} commands.
% Install: texlive-publishers or MiKTeX package manager.

\documentclass[preprint,12pt,3p]{elsarticle}
% Options:
%   preprint  — double-spaced, single column (use for submission)
%   review    — same as preprint but with line numbers
%   1p,2p,3p  — final 1/2/3-column layout (do NOT use for submission)

% ── Core packages ─────────────────────────────────────────────────────────────
\usepackage{amsmath,amssymb}
\usepackage{graphicx}
\usepackage{booktabs}
\usepackage{caption}
\usepackage{subcaption}
\usepackage{multirow}
\usepackage{siunitx}           % SI units
\usepackage{hyperref}
\usepackage{lineno}            % Line numbers for review
\modulolinenumbers[5]

% ── Hyperref ──────────────────────────────────────────────────────────────────
\hypersetup{colorlinks=true, linkcolor=blue, citecolor=blue, urlcolor=blue}

% ── SI units setup ────────────────────────────────────────────────────────────
\sisetup{separate-uncertainty=true, range-phrase=--}

% ══════════════════════════════════════════════════════════════════════════════
\journal{Engineering Structures}   % ← Change to target journal name

\begin{document}

\begin{frontmatter}

% ── Title ─────────────────────────────────────────────────────────────────────
\title{Insert Your Title Here: Concise, Specific, and Reflecting Main Contribution}

% ── Authors and Affiliations ──────────────────────────────────────────────────
\author[inst1]{First Author}
\author[inst2]{Second Author}
\author[inst1,inst3]{Corresponding Author\corref{cor1}}
\ead{corresponding@university.edu}

\cortext[cor1]{Corresponding author. Tel.: +XX-XXX-XXXX-XXXX.}

\affiliation[inst1]{organization={Department of Civil Engineering},
             city={City},
             country={Country}}

\affiliation[inst2]{organization={Department of Building Science},
             city={City},
             country={Country}}

\affiliation[inst3]{organization={Research Institute for Sustainable Infrastructure},
             city={City},
             country={Country}}

% ── Abstract ─────────────────────────────────────────────────────────────────
% Elsevier engineering journals: typically ≤200–250 words, unstructured
% Highlight: background (1–2 sentences), objectives (1 sentence),
%            methods (1–2 sentences), results (2–3 sentences with numbers),
%            conclusions (1 sentence)
\begin{abstract}
Write a concise abstract (200–250 words).
State the engineering problem and motivation in 1–2 sentences.
Define the specific objectives of this study.
Briefly describe the methodology (e.g., experimental program, FEM, LCA, ML model).
Present quantitative key results (e.g., \SI{23}{\percent} reduction in embodied carbon,
RMSE of \SI{0.045}{\kilo\watt\hour\per\meter\squared} for energy prediction).
Conclude with the significance and practical implications.
\end{abstract}

% ── Keywords ─────────────────────────────────────────────────────────────────
% Elsevier: 4–8 keywords, comma-separated; avoid words already in title
\begin{keyword}
keyword one \sep keyword two \sep keyword three \sep keyword four \sep keyword five
\end{keyword}

\end{frontmatter}

% ── Line numbers (activate for review submission) ─────────────────────────────
\linenumbers

% ══════════════════════════════════════════════════════════════════════════════
\section{Introduction}
% Recommended: 600–1000 words
% Structure: (1) broad context, (2) review of prior work with citations,
%            (3) identified research gap, (4) objectives of this study,
%            (5) paper organization overview

Provide engineering motivation and context~\cite{Author2023}.
Review prior work systematically, highlighting advances and remaining gaps~\cite{Smith2021,Jones2022}.

State the research gap and objectives clearly.

The remainder of this paper is organized as follows.
Section~\ref{sec:methods} describes \ldots

% ══════════════════════════════════════════════════════════════════════════════
\section{Methodology}
\label{sec:methods}
% Alternate section titles depending on study type:
%   "Experimental Program", "Numerical Framework", "Life-Cycle Assessment Framework",
%   "Machine Learning Model", "Energy Simulation Methodology"

\subsection{Study Scope and System Description}
Describe the system, material, structure, or process under investigation.
Report key parameters with SI units using \verb|\SI{}{}|:
e.g., floor area $= \SI{500}{\meter\squared}$,
concrete strength $f'_c = \SI{35}{\mega\pascal}$.

\subsection{Data Sources and Materials}
Describe data collection, material characterization, or input databases.
For LCA studies: declare functional unit, system boundaries,
LCI database (ecoinvent \cite{Wernet2016}, GaBi \cite{Sphera2024}),
and impact assessment method (ReCiPe 2016 \cite{Huijbregts2017}, IPCC 2021).
For ML studies: describe dataset origin, size ($n = $), features, and preprocessing.

\subsection{Analysis Method}
Describe numerical, simulation, experimental, or statistical approach.
Report software and version. For FEM: element types, mesh size, convergence criteria.
For building energy simulation: software (EnergyPlus v23.1 \cite{EnergyPlus2023}),
weather file (TMY3/EPW), HVAC system model.
For ML: architecture, training/validation/test split, hyperparameter tuning strategy.

\subsection{Validation}
Validate the method against experimental data, published benchmarks, or analytical solutions.
Report error metrics: $R^2$, RMSE (\si{\kilo\watt\hour\per\meter\squared}),
MAE, or coefficient of variation (CV).

% ══════════════════════════════════════════════════════════════════════════════
\section{Results}
\label{sec:results}
% Present findings without interpretation; save for Discussion
% Use subsections to organize by finding or by parameter

\subsection{Result Subsection One}
Present quantitative results.
Refer to Fig.~\ref{fig:example} and Table~\ref{tab:example}.

\begin{figure}[htbp]
  \centering
  \includegraphics[width=0.85\textwidth]{figures/figure_01.png}
  \caption{Figure caption. Describe axes, units, key trends.
           (a) Condition A; (b) Condition B.}
  \label{fig:example}
\end{figure}

\begin{table}[htbp]
  \centering
  \caption{Table title. Units stated in column headers.}
  \label{tab:example}
  \begin{tabular}{lcccc}
    \toprule
    Case & EUI (\si{\kilo\watt\hour\per\meter\squared}) & CO$_2$e (\si{\kilo\gram\per\meter\squared}) & $\Delta$EUI (\%) & $\Delta$CO$_2$e (\%) \\
    \midrule
    Baseline           & 120.4 & 58.2 & ---   & ---   \\
    Scenario A         &  98.1 & 47.5 & -18.5 & -18.4 \\
    Scenario B         &  85.3 & 41.3 & -29.2 & -29.0 \\
    Net-zero target    &  60.0 & 29.0 & -50.2 & -50.2 \\
    \bottomrule
  \end{tabular}
\end{table}

\subsection{Result Subsection Two}
Continue with additional results.

% ══════════════════════════════════════════════════════════════════════════════
\section{Discussion}
\label{sec:discussion}
% Interpret, compare, explain, and contextualize results
% Structure: (1) key findings explained, (2) comparison with literature,
%            (3) practical implications, (4) limitations and uncertainty

Interpret the results in light of the research objectives.
Compare findings with prior work~\cite{PriorWork2022}.
Discuss practical implications for engineering design, policy, or practice.
Acknowledge limitations (system boundary, dataset size, model assumptions).

% ══════════════════════════════════════════════════════════════════════════════
\section{Conclusions}
\label{sec:conclusions}
% Elsevier engineering: numbered list strongly preferred; concise and quantitative

The following conclusions are drawn from this study:
\begin{enumerate}
  \item First conclusion — quantitative result and its significance.
  \item Second conclusion — design implication or policy recommendation.
  \item Third conclusion — conditions of applicability or generalizability.
  \item Recommendation for future work.
\end{enumerate}

% ══════════════════════════════════════════════════════════════════════════════
% CRediT Author Contribution Statement (required by Elsevier since 2020)
\section*{CRediT authorship contribution statement}
\textbf{First Author:} Conceptualization, Methodology, Software, Formal analysis,
Writing -- original draft.
\textbf{Second Author:} Validation, Data curation, Visualization.
\textbf{Corresponding Author:} Conceptualization, Supervision, Funding acquisition,
Writing -- review \& editing.

% Declaration of Competing Interest
\section*{Declaration of competing interest}
The authors declare that they have no known competing financial interests or personal
relationships that could have appeared to influence the work reported in this paper.

% Data Availability Statement (required)
\section*{Data availability}
Data will be made available on request.
% OR: The data and code are publicly available at [DOI/URL].

% Acknowledgements
\section*{Acknowledgements}
This work was supported by [Funding Agency] under Grant [No.\ XXXXXX].
The authors thank [collaborators/facilities] for [contribution].

% ══════════════════════════════════════════════════════════════════════════════
% ELSEVIER REFERENCE FORMAT (numbered, sorted by appearance):
% [1] A. Smith, B. Jones, "Title of article," J. Clean. Prod. 280 (2021) 124320.
%     https://doi.org/10.1016/j.jclepro.2020.124320
%
% Example BibTeX entry:
% @article{Smith2021,
%   author  = {Smith, Alice and Jones, Bob},
%   title   = {Life-cycle carbon assessment of high-performance building envelopes},
%   journal = {J. Clean. Prod.},
%   year    = {2021},
%   volume  = {280},
%   pages   = {124320},
%   doi     = {10.1016/j.jclepro.2020.124320}
% }

\bibliographystyle{elsarticle-num}   % Numbered citations [1], [2]...
\bibliography{../references/references}

\end{document}
